% =====================================================================
%  26.5  The method boundary -- when a virtual tree is the WRONG tool.
%  Exports: sec:c30-boundary  def:c30-pick
%  WORKED examples only -- NO solvedproblem (in-section solved stays 2).
%  Discharges the forward-cites promised by Ch.25 (sec:c23-boundary hands the
%  "shrink a fixed tree" case here) and 26.3/26.4 (graph Steiner routed out).
%  Mirrors the ch23 05_boundary / ch18 05_boundary worked-only style.
% =====================================================================
\section{The Method Boundary: When a Virtual Tree Is the Wrong Tool}\label{sec:c30-boundary}

\begin{previouslybox}[Every engine ends by naming what it must hand away]
Chapter~25 closes its link-cut chapter by listing the specs an LCT must not touch
(\(\secref{c23-boundary}\)); Chapter~24 does the same for centroid decomposition.  A
virtual tree is the cheapest of the five tree engines to code and the most seductive to
over-apply: once you can see a ``marked set,'' you want to compress everything.  But the
compression only earns its keep when the answer is truly confined to the marked
set's induced subtree.  The mature reflex is not ``can I build \(A(S)\)?'' (you can, for
any \(S\)) but ``does \(A(S)\) \emph{lose} information the answer needs?''  This section
walks six specifications through the pick (\(\defref{c30-pick}\)); none has a solved
problem, because the whole lesson is \emph{not} reaching for a virtual tree.
\end{previouslybox}

\begin{topicintro}
A virtual tree is a \emph{lossy} projection: it keeps the marked nodes and their LCAs and
throws away every other real node.  That is free precisely when the discarded nodes
cannot change the answer -- which \(\thmref{c30-fold}\) guarantees for distances,
terminal counts and connected structures among the marked set.  The moment the
statement needs the discarded mass (a value summed over all nodes), a \emph{graph} where
terminals meet at new nodes, a topology that moves, or a range you never shrink, the
projection is wrong and a different engine is right.  Five such engines already exist in
this book; the section's job is to route to them by name.

\textbf{Prerequisites.} The whole of \(\secref{c30-construction}\)--\(\secref{c30-apps}\),
so you know what a virtual tree keeps before deciding it loses too much;
\(\secref{c30-recognition}\) (the same gate read as an acceptance test, here read as a
rejection test); Ch.~13 (\chref{c10}) tree DP, Ch.~15 (\chref{c20}) rerooting, Ch.~24
(\chref{c18}) centroid, Ch.~22 (\chref{c16}) HLD, Ch.~5 (\chref{c24}) Euler, Ch.~25
(\chref{c23}) link-cut.

\textbf{After this section you can}
\begin{itemize}[nosep]
  \item read a spec and say what the compression would \emph{throw away}; \hfill\lv{5}
  \item route ``every node matters'' to tree DP / rerooting; \hfill\lv{5}
  \item route graph terminals to bitmask Steiner and moving trees to link-cut. \hfill\lv{5}
\end{itemize}
\end{topicintro}

\begin{definition}{The virtual-tree tool pick}{c30-pick}
Ask three questions in order. \textbf{(Q1) Is the tree fixed and the answer confined to a
small marked set's induced subtree?} Yes \(\Rightarrow\) \emph{virtual tree}
(\(\secref{c30-construction}\)--\(\secref{c30-apps}\)); this is the only branch where the
compression is exact (\(\thmref{c30-fold}\)). \textbf{(Q2) If the answer needs \emph{every}
node's own value} (a global sum, a per-node count, a value that depends on unmarked
subtrees)\(\Rightarrow\) \emph{whole-tree DP} (\chref{c10}) or \emph{rerooting}
(\chref{c20}) -- no set to shrink. \textbf{(Q3) If the marked set is confined but the
structure is not a plain tree} \(\Rightarrow\) split by what breaks: terminals on a
\emph{graph} that meet at \emph{new} nodes \(\Rightarrow\) bitmask Steiner DP
(Dreyfus--Wagner \(O(3^k n)\)); a \emph{moving topology} \(\Rightarrow\) link-cut
(\chref{c23}); \emph{range/path} queries that never shrink with \(k\) \(\Rightarrow\) HLD
(\chref{c16}) or Euler+Fenwick (\chref{c24}); \emph{counting over all} marked pairs
without any per-set DP \(\Rightarrow\) possibly centroid (\chref{c18}) or Mo-on-trees
(\chref{c21}) instead.  Only Q1's ``yes'' is a virtual tree.
\end{definition}

\begin{comparisonbox}[Same question, five other homes: what each spec actually wants]
\setlength{\tabcolsep}{3pt}\renewcommand{\arraystretch}{1.2}%
\begin{tabular}{@{}p{0.30\linewidth}p{0.31\linewidth}p{0.31\linewidth}@{}}
\toprule
\textbf{Signal in the spec} & \textbf{Why \(A(S)\) fails} & \textbf{Right engine}\\
\midrule
answer is a sum over all \(n\) nodes & unmarked nodes carry the value & tree DP / reroot \(\chref{c10},\chref{c20}\)\\
terminals live on a \emph{graph} & connectors may be new nodes & bitmask Steiner DP\\
edges are linked / cut online & the tree is not static & link-cut \(\chref{c23}\)\\
fixed path/subtree \emph{ranges} & nothing shrinks with \(k\) & HLD / Euler \(\chref{c16},\chref{c24}\)\\
count over \emph{all} pairs, no DP & centroid/Mo already linear & centroid \(\chref{c18}\), Mo \(\chref{c21}\)\\
\bottomrule
\end{tabular}
\end{comparisonbox}

\subsection*{Six specs, read and routed}

\begin{itemize}[leftmargin=1.4em]
  \item \textbf{``Fixed tree; each node has a value; for every node report the sum of
        its subtree, or a rerooted total.''} The answer is defined on \emph{all} \(n\)
        nodes, not a marked set -- there is nothing to compress. \(\Rightarrow\) \emph{tree
        DP} (\chref{c10}) / \emph{rerooting} (\chref{c20}). \textbf{Do not.}
  \item \textbf{``A weighted \emph{graph}; connect \(k\le18\) terminals at minimum cost.''}
        On a graph the cheapest Steiner node can be a non-terminal not present in any
        \(A(S)\). \(\Rightarrow\) \emph{bitmask Steiner DP} (Dreyfus--Wagner). \textbf{Do
        not} -- this is the mixed-set CF 152E contrast (\(\secref{c30-apps}\)'s insight box).
  \item \textbf{``A tree whose edges are linked and cut between queries; ask a path
        aggregate online.''} The topology moves, so \(\tin/\tout/\lca\) are invalid after
        the first edit and there is no static order to sort by. \(\Rightarrow\)
        \emph{link-cut tree} (\chref{c23}). \textbf{Do not.}
  \item \textbf{``Fixed tree; many queries, each a whole path or whole subtree range
        sum/min with updates; nothing names a small marked set.''} Every query touches an
        \(O(n)\)-sized range, so compressing per query loses nothing to gain.
        \(\Rightarrow\) \emph{HLD} (\chref{c16}) or \emph{Euler+Fenwick} (\chref{c24}).
        \textbf{Prefer those.}
  \item \textbf{``Fixed tree; count pairs \((u,v)\) with a distance property over the
        \emph{whole} tree.''} There is no marked set -- it is all pairs -- so this is a
        distance-distribution job. \(\Rightarrow\) \emph{centroid decomposition}
        (\chref{c18}); \(\sum\) per virtual edge only helps a \emph{named} set, not the
        global count. \textbf{Do not.}
  \item \textbf{``One fixed tree; an offline batch of node-set queries with a slowly-updated
        statistic and no inverse.''} The sets overlap heavily and are answered offline.
        \(\Rightarrow\) \emph{Mo's algorithm on trees} (\chref{c21}) -- compressing each set
        independently re-does work Mo amortises away. \textbf{Prefer Mo} when the add/remove
        of a single node is cheap but there is no fold to run.
\end{itemize}

\begin{insightbox}[The honest cost of the compression]
A virtual tree is cheap to build and \emph{exact} only inside its promise: distances and
connected structures among the marked set are preserved (\(\thmref{c30-fold}\)), but
anything about the unmarked nodes is silently dropped.  The failure mode is not a wrong
answer you can see -- it is a right-looking number computed on the wrong (too small)
tree.  Before trusting \(A(S)\), re-read the statement and point at one unmarked node:
if it could ever change the answer, you do not have a virtual-tree problem, and Q2/Q3
hand you to a full-tree engine.
\end{insightbox}

\begin{warningbox}[The four ways students reach for a virtual tree too early]
\begin{enumerate}[nosep,leftmargin=1.6em]
  \item ``the problem names special vertices'' but the value is summed over \emph{all}
        nodes -- that is tree DP, not a compression (\chref{c10}).
  \item ``connect these terminals'' on a \emph{graph} -- terminals may join through new
        nodes, so \(A(S)\) has nowhere to put the connector; bitmask Steiner DP.
  \item ``a marked set per query'' where the tree is being \emph{edited} -- no static
        \(\tin\) order survives a cut; link-cut (\chref{c23}).
  \item ``over all pairs of marked nodes'' when the marked set \emph{is} the whole tree
        (\(k=n\)) -- \(A(S)=T\), the compression buys nothing; a plain fold or centroid.
\end{enumerate}
Each wastes a virtual tree on a job a full-tree or different engine does correctly.
\end{warningbox}

\begin{hardeningbox}
\begin{itemize}[nosep,leftmargin=1.4em]
  \item \textbf{Point at a discarded node.} Name one node not in \(A(S)\) and ask if the
        answer could depend on it. If yes, abandon the compression before coding it.
  \item \textbf{Check \(k\ll n\) really holds.} If the marked set is a constant fraction
        of \(n\), \(A(S)=O(n)\) and you paid \(\bigO{k\log k}\) to build a tree you then
        walk in full -- the plain \(O(n)\) DP is faster.
  \item \textbf{``Connect the terminals'' is two problems.} On a tree it is free
        (\(\secref{c30-apps}\), the expanded \(A(S)\)); on a graph it is NP-hard (bitmask
        DP).  The single word \emph{graph} flips the engine.
\end{itemize}
\end{hardeningbox}

\begin{recallbox}[26.5 in one breath]
\emph{A virtual tree is a lossy projection that is exact only for distances and connected
structures among a marked set on a fixed tree. Every node matters \(\to\) tree DP /
reroot; terminals on a graph meet at new nodes \(\to\) bitmask Steiner; the topology
moves \(\to\) link-cut (\(\chref{c23}\)); ranges never shrink \(\to\) HLD / Euler; all
pairs with no per-set DP \(\to\) centroid or Mo. Three questions decide it: is the tree
fixed, is the answer confined to a small marked set, and would dropping the unmarked
mass lose anything (\(\defref{c30-pick}\)).}
\selftest{
  \item A spec has a marked set but sums a value over every node -- which engine? \emph{(Tree DP / rerooting, \chref{c10}/\chref{c20}; the compression drops nodes the answer needs.)}
  \item ``Connect \(k\) terminals'' -- when is it \emph{not} a virtual tree? \emph{(When they sit on a graph and may join at new non-terminal nodes: bitmask Steiner DP, \(\secref{c30-boundary}\).)}
  \item State the three questions of the pick. \emph{(Fixed tree? Answer confined to a small marked set's induced subtree? Does dropping unmarked nodes lose information? \(\defref{c30-pick}\).)}
}
\end{recallbox}
